% !TEX root = ../LAFROX-Formulario-T-17.tex
\chapter*{Formulario T-17: Protocolo de aceptación}
\addcontentsline{toc}{chapter}{Formulario T-17: Protocolo de aceptación}
\markboth{Formulario T-17: Protocolo de aceptación}{}
\label{cap:t17}

Este formulario acompaña al Subdocumento 9, Plan de calidad (SD9). Contiene el protocolo de aceptación de cada hito y del producto final que exige el Formulario T-17 de las Bases Administrativas: entregables, criterios de aceptación objetivos, evidencia requerida, plazos de revisión, procedimiento de observaciones y acta de conformidad (Bases Administrativas, Formulario T-17, p. 64; Bases Técnicas Transversales, RT-20.08, p. 35).

El protocolo recibe los entregables que el control de calidad ya verificó con el plan de pruebas del Formulario T-13. Verificar le corresponde a LafroX; aceptar le corresponde al CLIENTE, por medio de su Contraparte Técnica (Project Management Institute [PMI], 2017, p. 305).

Mes 1 del contrato = febrero de 2027. Paquetes de trabajo del Formulario T-14; fechas y holguras del T-15, Tabla 8 y secciones 5.3 a 5.5. La sección 4.1 de este formulario reproduce las fechas del SD7 y distingue entrega, acta y límite contractual.

\section*{1. Base contractual}
\addcontentsline{toc}{section}{1. Base contractual}

Las reglas siguientes vinculan cada decisión del protocolo con su fundamento y con la evidencia que debe acompañar el hito.

\begin{tablasimple}{>{\hsize=1.733\hsize\linewidth=\hsize}Y >{\hsize=0.318\hsize\linewidth=\hsize}Y >{\hsize=0.949\hsize\linewidth=\hsize}Y}
  {\cab{Regla} & \cab{Fuente} & \cab{Aplicación en este protocolo}}
Un hito se cumple solo con el acta firmada por la Contraparte Técnica & Art. 18.1 (BA, p. 13) & Ningún hito se informa como cumplido sin acta \\
Todo entregable incluye el artefacto, la evidencia objetiva de su verificación, la trazabilidad a los requisitos que satisface y el registro de las observaciones previas resueltas & Art. 18.2 (BA, p. 13) & Contenido obligatorio del expediente (sección 3) \\
10 días hábiles de revisión del CLIENTE y 10 días hábiles de subsanación; repetir una observación de la misma naturaleza es atraso imputable & Art. 18.3 (BA, p. 13) & Procedimiento de la sección 2 \\
La aceptación no libera de responsabilidad por defectos posteriores & Art. 18.4 (BA, p. 13) & Los defectos posteriores se tratan como incidentes de la Operación \\
Seis condiciones copulativas para cerrar una marcha blanca & Art. 17.3 (BA, p. 13) & Acta de cierre de cada marcha blanca (paquetes 4.2.1 y 4.3.1), que habilita el paso a producción previo al H7 y al H12 \\
Definición de terminado acordada con el CLIENTE & RT-20.07 (BTT, p. 35) & Sección 3.2 \\
Protocolo formalizado conforme al T-17, con criterios objetivos y verificables & RT-20.08 (BTT, p. 35) & Todo este formulario \\
\end{tablasimple}
\fuentesd{Fuente: (Bases Administrativas, arts. 17.3, 18 y 70.2, pp. 13, 37); (Bases Técnicas Transversales, RT-20.07 y RT-20.08, p. 35).}

\section*{2. Procedimiento de aceptación}
\addcontentsline{toc}{section}{2. Procedimiento de aceptación}

La aceptación exige entrega formal, revisión y pronunciamiento del CLIENTE. Las revisiones preliminares ayudan a preparar el expediente, pero no sustituyen esos pasos.

\subsection*{2.1 Pasos}
\addcontentsline{toc}{subsection}{2.1 Pasos}

El procedimiento distingue conformidad sin observaciones y subsanación. Ningún hito se declara cumplido antes de su acta.

\begin{tablasimple}{>{\hsize=0.89\hsize\linewidth=\hsize}Y >{\hsize=0.592\hsize\linewidth=\hsize}Y >{\hsize=1.178\hsize\linewidth=\hsize}Y >{\hsize=1.34\hsize\linewidth=\hsize}Y}
  {\cab{Paso} & \cab{Responsable} & \cab{Plazo} & \cab{Resultado}}
1. Verificación interna del expediente con la lista de la sección 3.3 & Líder de Calidad y dueño del entregable & Antes de la entrega & Expediente completo y lista firmada internamente \\
2. Entrega formal del expediente & Jefe de Proyecto & En la fecha programada de la sección 4, no después de su límite; el límite reserva 10 días hábiles de revisión del CLIENTE hasta el cierre del mes & Constancia de recepción \\
3. Revisión del CLIENTE & Contraparte Técnica & 10 días hábiles (BA, Art. 18.3, p. 13) & Conformidad u observaciones por escrito \\
4a. Sin observaciones & Contraparte Técnica & — & Acta de conformidad (sección 6) \\
4b. Con observaciones & LafroX & 10 días hábiles contractuales; la holgura del hito no reduce este plazo ni garantiza que quepa una subsanación & Expediente subsanado y registro de cada observación resuelta \\
5. Segunda revisión & Contraparte Técnica & Sin presumir revisión instantánea; fecha acordada y acta requerida & Acta, o atraso imputable si se repite una observación de la misma naturaleza \\
\end{tablasimple}
\fuentesd{Fuente: (Bases Administrativas, arts. 18.1 a 18.4, p. 13); (Formulario T-15, sección 5.5); (Formulario T-13, sección 6).}

\subsection*{2.2 Tratamiento de las observaciones}
\addcontentsline{toc}{subsection}{2.2 Tratamiento de las observaciones}

Cada observación se registra con su identificador, el criterio de aceptación al que se refiere, la evidencia que la sustenta, su clasificación, su causa y su responsable.

\begin{tablasimple}{>{\hsize=0.333\hsize\linewidth=\hsize}Y >{\hsize=1.464\hsize\linewidth=\hsize}Y >{\hsize=1.203\hsize\linewidth=\hsize}Y}
  {\cab{Clasificación} & \cab{Definición} & \cab{Efecto en el acta}}
Bloqueante & El entregable no cumple un criterio de aceptación de la sección 4 o 5 & Impide el acta hasta su subsanación \\
No bloqueante & Mejora de forma o de contenido que no afecta un criterio de aceptación & Se acepta con el acta; se cierra en un plazo acordado y registrado en ella \\
\end{tablasimple}
\fuentesd{Fuente: (Bases Administrativas, art. 18.3, p. 13); (Formulario T-13, sección 6); los criterios de hito de las secciones 4 y 5 de este protocolo.}

La clasificación la propone la Contraparte Técnica y queda en el registro. Toda observación se analiza por su causa y se agrega a la lista de verificación de la sección 3.3 para que no se repita en otro entregable.

\subsection*{2.3 Hitos con reserva menor que la subsanación}
\addcontentsline{toc}{subsection}{2.3 Hitos con reserva menor que la subsanación}

H1, H7 y H12 tienen holguras contractuales publicadas de cinco, dos y tres días hábiles, menores que los diez de subsanación. Sus expedientes se revisan con la Contraparte Técnica antes de la entrega formal. La versión preliminar del expediente H1 se prepara del 1 al 4 de marzo de 2027, se revisa el 5 y se corrige y cierra el 8 y el 9 dentro de las 80 HH de 1.5.1. En H7 y H12, la preparación del expediente avanza durante las últimas semanas de marcha blanca, sin anticipar su aceptación ni sustituir las cuatro semanas finales de evidencia. H8 dispone de catorce días de holgura. Se confirma asistencia del CLIENTE. Las revisiones preliminares no sustituyen sus diez días de revisión formal, los diez de subsanación ni una segunda revisión. Si hay observaciones, se registra y escala su impacto sin garantizar que el acta quepa dentro del mes.

\subsection*{2.4 Entrega por incrementos}
\addcontentsline{toc}{subsection}{2.4 Entrega por incrementos}

Los entregables que admiten revisión por partes se presentan por incrementos. El H4 y el H9 se presentan con el informe de QA de cada módulo al terminarlo, y la revisión final cubre solo el último incremento (Formulario T-15, sección 5.5).

\section*{3. Expediente de aceptación}
\addcontentsline{toc}{section}{3. Expediente de aceptación}

El expediente permite reconstruir qué se entrega, qué se verificó y qué observaciones se resolvieron; se comprueba antes de su recepción formal.

\subsection*{3.1 Contenido mínimo (Art. 18.2)}
\addcontentsline{toc}{subsection}{3.1 Contenido mínimo (Art. 18.2)}

Los cuatro componentes siguientes acompañan toda presentación. La evidencia debe corresponder a la misma versión del artefacto entregado.

\begin{tablasimple}{>{\hsize=0.361\hsize\linewidth=\hsize}Y >{\hsize=1.639\hsize\linewidth=\hsize}Y}
  {\cab{Componente} & \cab{Contenido}}
Artefacto & Documento, software desplegado, informe o configuración, en la versión que se acepta, identificado por versión o por el digest de la imagen \\
Evidencia objetiva & Informes de las pruebas del T-13 que aplican al hito, con resultado medido frente a su umbral \\
Trazabilidad & Extracto de la matriz de trazabilidad: requisitos del T-12 que el entregable satisface, con su caso de prueba y su resultado \\
Observaciones previas & Registro de las observaciones de entregas anteriores, con su resolución \\
\end{tablasimple}
\fuentesd{Fuente: (Bases Administrativas, art. 18.2, p. 13); (Bases Técnicas Transversales, RT-04.04, p. 10); (Formulario T-12, matriz de trazabilidad).}

\subsection*{3.2 Definición de terminado (RT-20.07)}
\addcontentsline{toc}{subsection}{3.2 Definición de terminado (RT-20.07)}

Un entregable de software está terminado cuando cumple las cinco condiciones siguientes. Se acuerdan con la Contraparte Técnica en el H1 y no se modifican sin acta.

\begin{tablasimple}{>{\hsize=0.31\hsize\linewidth=\hsize}Y >{\hsize=1.69\hsize\linewidth=\hsize}Y}
  {\cab{Dimensión} & \cab{Condición}}
Código & Fusionado a la rama principal con aprobación de un par; umbrales de código del T-13, sección 6.1, cumplidos \\
Pruebas & Casos del requisito aprobados; regresión completa aprobada; 0 defectos críticos o altos abiertos \\
Documentación & Contratos de interfaz, manual de la funcionalidad y matriz de trazabilidad actualizados \\
Seguridad & 0 hallazgos críticos o altos abiertos en SAST, SCA, DAST e imágenes \\
Despliegue & Imagen firmada promovida sin recompilar desde QA (RT-04.08); despliegue y reversión probados (BTT, p. 11) \\
\end{tablasimple}
\fuentesd{Fuente: (Bases Técnicas Transversales, RT-20.07 y RT-04.08, pp. 11, 35); (Formulario T-13, sección 6.1); (Formulario T-10, apartado DevSecOps).}

\subsection*{3.3 Lista de verificación antes de la entrega}
\addcontentsline{toc}{subsection}{3.3 Lista de verificación antes de la entrega}

El Líder de Calidad aplica esta lista al expediente y registra el resultado antes de presentarlo a la Contraparte Técnica.

\begin{tablasimple}{>{\hsize=0.54\hsize\linewidth=\hsize}Y >{\hsize=1.46\hsize\linewidth=\hsize}Y}
  {\cab{Comprobación} & \cab{Criterio}}
Correspondencia con las Bases y con el alcance ofertado & Cada criterio del hito tiene su evidencia en el expediente \\
Cifras, unidades y cálculos & Mismos valores en el artefacto, en sus informes y en los formularios relacionados \\
Trazabilidad & Requisito, componente, paquete, prueba y criterio vinculados sin vacíos \\
Estado real & El expediente distingue compromiso, prueba prevista, resultado obtenido y aceptación; no atribuye aprobaciones sin su acta \\
Observaciones previas & Todas las observaciones de entregas anteriores resueltas y registradas \\
\end{tablasimple}
\fuentesd{Fuente: (Bases Administrativas, arts. 18.2 y 18.3, p. 13); (Bases Técnicas Transversales, RT-20.08, p. 35); los criterios de hito y expediente de las secciones 3 a 5 de este protocolo.}

\section*{4. Protocolo por hito}
\addcontentsline{toc}{section}{4. Protocolo por hito}

Se distinguen los hitos documentales de la entrada al ambiente productivo y de la aceptación de las marchas blancas.

\subsection*{4.1 Fechas de entrega}
\addcontentsline{toc}{subsection}{4.1 Fechas de entrega}

La tabla distingue entrega, acta calculada y límite contractual del acta. La holgura se mide después de los diez días hábiles de revisión del CLIENTE; no es otro plazo de revisión ni una reserva de personal.

\begin{tablasimple}{Y Y Y Y Y}
  {\cab{Hito (mes)} & \cab{Entrega calculada} & \cab{Acta calculada} & \cab{Límite del acta} & \cab{Holgura (días hábiles)}}
H1 (2) & 09-03-2027 & 23-03-2027 & 31-03-2027 & 5 \\
H2 (4) & 29-04-2027 & 13-05-2027 & 31-05-2027 & 11 \\
H3 (6) & 18-06-2027 & 06-07-2027 & 30-07-2027 & 17 \\
H4 (10) & 22-10-2027 & 08-11-2027 & 30-11-2027 & 16 \\
H5 (12) & 01-12-2027 & 03-01-2028 & 31-01-2028 & 20 \\
H8 (14) & 28-02-2028 & 13-03-2028 & 31-03-2028 & 14 \\
H7 (16) & 15-05-2028 & 29-05-2028 & 31-05-2028 & 2 \\
H9 (17) & 18-05-2028 & 01-06-2028 & 30-06-2028 & 19 \\
H10 (18) & 28-06-2028 & 12-07-2028 & 31-07-2028 & 13 \\
H12 (21) & 11-10-2028 & 25-10-2028 & 31-10-2028 & 3 \\
\end{tablasimple}
\fuentesd{Fuente: (Bases Administrativas, art. 18.3 y Formulario E-25, pp. 13, 74--75); (Formulario T-15, secciones 5.2 y 5.5); las reglas de calendario y feriados explicitadas en esta sección 4.1.}

El calendario administrativo, las fechas y las holguras se toman del T-15, Tabla 8. Se distingue entrega de acta: H5 se entrega en diciembre y su acta se calcula en enero, mes 12. La holgura se mide desde el acta calculada hasta su límite contractual, no desde la entrega hasta un límite documental. H6/H11 inician marcha blanca y H7/H12 exigen las condiciones de aceptación de la sección 4.3. La cobertura de Operación de ambos alcances comienza desde el primer día del mes 21, conforme a T-18; el acta H12 se presenta con la fecha calculada del T-15 y no se usa para diferir esa cobertura.

\subsection*{4.2 Hitos de entregables (H1 a H5 y H8 a H10)}
\addcontentsline{toc}{subsection}{4.2 Hitos de entregables (H1 a H5 y H8 a H10)}

Cada hito documental se acepta por su artefacto y por la evidencia asociada; la presentación de un informe no demuestra por sí sola el cumplimiento.

\begin{tablasimple}{>{\hsize=0.236\hsize\linewidth=\hsize}Y >{\hsize=1.288\hsize\linewidth=\hsize}Y >{\hsize=1.288\hsize\linewidth=\hsize}Y >{\hsize=1.188\hsize\linewidth=\hsize}Y}
  {\cab{Hito} & \cab{Entregable y paquete T-14} & \cab{Criterios de aceptación} & \cab{Evidencia}}
H1 (mes 2) & Línea base de alcance (1.2.1) y matriz de trazabilidad (1.2.4); plan de calidad y de pruebas (1.5.1) & Ningún requisito ofertado del T-12 queda sin paquete ni sin prueba prevista; cada tipo de prueba de la BTT 20.1 tiene su ambiente, calendario y criterio de salida; definición de terminado acordada (BTT, pp. 34--35) & Matriz; plan; acta de la revisión de la versión preliminar \\
H2 (mes 4) & Documento de arquitectura, plan de seguridad y modelo de datos (2.4.1) & Arquitectura con las cinco vistas requeridas por RT-02.03 y documentadas según ISO/IEC/IEEE 42010:2022 (ISO, 2022) y registro de decisiones; modelo de datos con su diccionario; el diseño incorpora los hallazgos de las pruebas de usabilidad de 2.6.2 (RT-13.03) (BTT, p. 26) & Documentos; informe de usabilidad de 2.6.2 como antecedente \\
H3 (mes 6) & Infraestructura híbrida y cinco ambientes con observabilidad (3.1.1, 3.1.2, 3.1.3 y 3.1.5); cadena de integración con puertas activas (1.5.2) & Ambientes habilitados (RT-04.01); Preproducción equivalente a Producción (RT-04.02); una versión con cobertura menor al umbral o con hallazgo crítico queda bloqueada automáticamente (BTT, p. 10) & Inventario de ambientes; ejecución de la cadena que demuestra un bloqueo \\
H4 (mes 10) & Software de la Etapa 1 con QA superado (3.8.1) & 100 \% de los flujos de E1 sin error, según Anexo 7.E; portales y EDI se aceptan en E2; 0 regresiones; cobertura ≥ 70 \% en lógica de negocio, ≥ 80 \% global y ≥ 80 \% de líneas modificadas & Informes de QA por módulo e informe de integración \\
H5 (mes 12) & Certificación de la Etapa 1 (3.8.7) & Pruebas 3.8.2 a 3.8.6 y 3.8.8 aprobadas con los criterios del T-13, secciones 6 y 7, incluidas las actas de los protocolos del SD4, Anexo 4-V (T-13, sección 7.5); 0 defectos críticos o altos abiertos & Informes de aceptación, perfil operacional, carga y resiliencia, recuperación, intrusión y despliegue; casos firmados \\
H8 (mes 14) & Línea base de alcance de la Etapa 2 (1.2.5) y su diseño detallado (2.4.2) & Matriz de trazabilidad actualizada para la Etapa 2; diseño consistente con la arquitectura del H2 & Matriz; diseño; acta de la revisión de la versión preliminar \\
H9 (mes 17) & Software de la Etapa 2 con QA superado (3.9.1) & 15 de 15 interfaces sin error con el alcance de E2 integrado, según Anexo 7.E; 0 regresiones en la Etapa 1 en producción & Informes de QA e informe de regresión de la Etapa 1 \\
H10 (mes 18) & Certificación de la Etapa 2 y cierre del desarrollo (3.9.6) & Pruebas 3.9.2 a 3.9.5 y 3.9.7 aprobadas, incluidos el perfil operacional del alcance cambiado y los protocolos del SD4, Anexo 4-V, que aplican a la Etapa 2; 0 defectos críticos o altos abiertos & Informes y casos firmados \\
\end{tablasimple}
\fuentesd{Fuente: (Bases Administrativas, Formulario E-25, pp. 74--75); (Bases Técnicas Transversales, sección 20.1, pp. 34--35); (Formulario T-14, paquetes indicados en cada fila); (Formulario T-13, secciones 6 y 7).}

\subsection*{4.3 Hitos de operación supervisada (H6, H7, H11 y H12)}
\addcontentsline{toc}{subsection}{4.3 Hitos de operación supervisada (H6, H7, H11 y H12)}

H6 y H11 habilitan la operación supervisada con G5; H7 y H12 aceptan el cierre contractual de la marcha blanca con las seis condiciones copulativas.

\begin{tablasimple}{>{\hsize=0.997\hsize\linewidth=\hsize}Y >{\hsize=0.252\hsize\linewidth=\hsize}Y >{\hsize=1.376\hsize\linewidth=\hsize}Y >{\hsize=1.375\hsize\linewidth=\hsize}Y}
  {\cab{Hito} & \cab{Paquete T-14} & \cab{Criterios de aceptación} & \cab{Evidencia}}
H6 (mes 13): entrada de la Etapa 1 a Producción (puerta G5) e inicio de su marcha blanca & 4.2.1 & Puerta G5 superada (SD9, sección 9.2.1); H5 aceptado; procedimiento de reversión ensayado (4.1.2); usuarios de la primera ola capacitados; conciliación de datos migrados sin diferencias no explicadas (3.7.5) & Acta del H5; informe del ensayo de reversión; registro de capacitación; acta de corte \\
H7 (mes 16): aceptación de la Etapa 1 en producción & 4.2.3 & Acta de cierre de la marcha blanca (4.2.1), firmada al término del mes 15 con las seis condiciones del Art. 17.3, medidas con los indicadores diarios del SD7, Tabla 7.9, durante las cuatro últimas semanas; paso a producción del 5 de mayo de 2028; resultados del caso de la sección 5 previstos para este hito (BA, p. 13) & Informes diarios de indicadores y de conciliación; registro de certificación de usuarios (7.3.1); informes de los resultados del caso \\
H11 (mes 19): entrada de la Etapa 2 a Producción (puerta G5) e inicio de su marcha blanca & 4.3.1 & Puerta G5 superada (SD9, sección 9.2.1); H10 aceptado; reversión ensayada; usuarios capacitados; fuente única de datos entre ambas etapas & Acta del H10; informe de reversión; registro de capacitación \\
H12 (mes 21): aceptación final del proyecto de implementación & 4.3.3 y 4.3.4 & Acta de cierre de la marcha blanca de la Etapa 2 (4.3.1), firmada al término del mes 20 con las seis condiciones del Art. 17.3; paso a producción del 5 de octubre de 2028; sin doble digitación; la Etapa 1 no se degradó (BA, Art. 17.2, p. 13); garantía de correcto funcionamiento constituida & Informes diarios; comparación de los indicadores de la Etapa 1 antes y durante la Etapa 2; certificación de usuarios (7.3.2); garantía \\
\end{tablasimple}
\fuentesd{Fuente: (Bases Administrativas, arts. 17.2 y 17.3 y Formulario E-25, pp. 13, 74--75); (Formulario T-14, cuenta 4); (Formulario T-13, secciones 6 y 9.2).}

Los indicadores del Art. 17.3 son los de la Tabla 7.9 del SD7: (BA, p. 13)

\begin{tablasimple}{>{\hsize=0.861\hsize\linewidth=\hsize}Y >{\hsize=1.345\hsize\linewidth=\hsize}Y >{\hsize=0.794\hsize\linewidth=\hsize}Y}
  {\cab{Condición del Art. 17.3} (BA, p. 13) & \cab{Indicador} & \cab{Umbral}}
1. Sin incidentes críticos ni altos abiertos & Incidentes críticos y altos abiertos atribuibles a la solución & 0 \\
2. Volumen real comprometido & Transacciones de la ola registradas en la solución & 100 \% del volumen real durante las cuatro últimas semanas \\
3. Disponibilidad y tiempo de respuesta & Indisponibilidad entre 05:30 y 07:00; disponibilidad de la transacción crítica; p95 & 0 minutos; ≥ 99,9 \%; preparación ≤ 1 s, entrega ≤ 2 s, preventa ≤ 1,5 s, stock y crédito ≤ 2 s \\
4. Conciliación & Diferencias sin explicar & 0 \\
5. Personal capacitado y certificado & Usuarios de la ola certificados & 100 \% \\
6. Acta & Acta de la Contraparte Técnica & Firmada \\
\end{tablasimple}
\fuentesd{Fuente: (Bases Administrativas, art. 17.3, p. 13); (Subdocumento 7, Tabla 7.9).}

Ninguna condición compensa a otra. Si una no se cumple al término de la marcha blanca, esta se extiende a costo de LafroX sin desplazar las fechas siguientes (BA, Art. 17.3, p. 13).

\section*{5. Aceptación de los resultados del caso (producto final)}
\addcontentsline{toc}{section}{5. Aceptación de los resultados del caso (producto final)}

El producto final se juzga por los dieciséis resultados del Caso, cap. 18 (Bases Técnicas del caso, cap. 18, pp. 34--35). La aceptación técnica de las certificaciones (H5 y H10) no acepta por sí sola estos resultados: cada uno se acepta con su medición en operación real, en el momento que fija el SD7, Tabla 7.11. Las metas son las del SD3, Anexo 3.J.

\begin{tablasimple}{>{\hsize=1.041\hsize\linewidth=\hsize}Y >{\hsize=0.936\hsize\linewidth=\hsize}Y >{\hsize=1.378\hsize\linewidth=\hsize}Y >{\hsize=0.645\hsize\linewidth=\hsize}Y}
  {\cab{Resultado} & \cab{Meta} & \cab{Medición} & \cab{Momento de aceptación}}
R18-01 Clientes afectados por un lote & < 2 horas & Simulacro de retiro con un lote real, cronometrado & Marcha blanca E1 \ensuremath{\rightarrow} H7 \\
R18-02 Recepciones con lote & 100 \% & Recepciones con lote / recepciones que lo requieren, cuatro semanas & Marcha blanca E1 \ensuremath{\rightarrow} H7 \\
R18-03 Registro continuo de temperatura & 100 \% de cámaras y vehículos con frío & Auditoría de series sin brechas y consulta por el CLIENTE & Marcha blanca E1 \ensuremath{\rightarrow} H7 \\
R18-04 OTIF medido de una sola forma & 90 \% (mes 15); 93 \% (mes 19); 95 \% (mes 32) & Medición diaria; revisión mensual en comité & Por tramo \\
R18-05 Stock y crédito en la preventa & 100 \% de pedidos con stock y crédito visibles & Registro de consultas por pedido & Marcha blanca E1 \ensuremath{\rightarrow} H7 \\
R18-06 Ningún pedido perdido o duplicado por falta de señal & 0 pérdidas y 0 duplicados & Prueba de 14 h sin señal y conciliación diaria & Marcha blanca E1 \ensuremath{\rightarrow} H7 \\
R18-07 Ruta del día siguiente & < 20 minutos & Medición diaria del tiempo de generación & Marcha blanca E1 \ensuremath{\rightarrow} H7 \\
R18-08 Prueba de entrega digital el mismo día & 100 \% de entregas & Entregas con prueba digital y constancia de recepción o consulta por el cliente & Marcha blanca E1 \ensuremath{\rightarrow} H7 \\
R18-09 Guías extraviadas o ilegibles & 0 \% & Localización y lectura de todas las guías del período & Marcha blanca E1 \ensuremath{\rightarrow} H7 \\
R18-10 Rendición cuadrada el mismo día & 100 \% de diferencias con causal el mismo día & Cierre diario por camión & Marcha blanca E1 \ensuremath{\rightarrow} H7 \\
R18-11 Costo de servir por cliente y entrega & 100 \% de entregas con costo calculado & Reconciliación con la contabilidad del mes & Marcha blanca E2 \ensuremath{\rightarrow} H12 \\
R18-12 Control de envases retornables & Pérdida anual ≤ 7 \% del parque & Saldo conciliado por cliente y transportista & Provisional en H7; final tras 12 meses de operación \\
R18-13 Pedidos de cadenas por EDI & 100 \% de los pedidos de cadenas certificadas; aviso de despacho en ≤ 2 minutos & Pedidos EDI / pedidos de las cadenas comprometidas & Marcha blanca E2 \ensuremath{\rightarrow} H12 \\
R18-14 Ocupación de camiones & Cada salida ≥ 60 \% de la capacidad útil; promedio > 68 \% & Ocupación de cada salida durante cuatro semanas & Marcha blanca E1 \ensuremath{\rightarrow} H7 \\
R18-15 Causa de cada faltante & 100 \% con causa registrada al producirse & Trazabilidad del evento, incluso sin conexión & Marcha blanca E1 \ensuremath{\rightarrow} H7 \\
R18-16 Operación sin el planificador & OTIF sin caída frente al período con el planificador & Dos semanas sin el planificador frente a dos semanas comparables & Marcha blanca E1 \ensuremath{\rightarrow} H7 \\
\end{tablasimple}
\fuentesd{Fuente: (Bases Técnicas del caso, cap. 18, pp. 34--35); (Subdocumento 3, anexo 3.J); (Formulario T-14, cuentas 4 y 8); (Formulario T-13, sección 7.4).}

\section*{6. Acta de conformidad}
\addcontentsline{toc}{section}{6. Acta de conformidad}

El acta se emite por hito y contiene lo siguiente. Un acta sin alguno de estos campos no acredita el hito.

\begin{tablasimple}{>{\hsize=0.372\hsize\linewidth=\hsize}Y >{\hsize=1.628\hsize\linewidth=\hsize}Y}
  {\cab{Campo} & \cab{Contenido}}
Identificación & Hito, mes contractual, versión o digest del artefacto aceptado \\
Criterios & Cada criterio de aceptación del hito con su resultado: cumplido o no cumplido \\
Evidencia & Referencia a cada informe del expediente \\
Observaciones & Observaciones no bloqueantes por cerrar, con responsable y plazo de cierre \\
Pronunciamiento & Conforme, con fecha \\
Firmas & Contraparte Técnica del CLIENTE y Jefe de Proyecto de LafroX \\
\end{tablasimple}
\fuentesd{Fuente: (Bases Administrativas, arts. 18.1 y 18.2, p. 13); (Bases Técnicas Transversales, RT-20.08, p. 35).}

\section*{7. Roles}
\addcontentsline{toc}{section}{7. Roles}

La Contraparte Técnica acepta; LafroX prepara, verifica y subsana. La tabla delimita esas responsabilidades sin transferir la aceptación al equipo de pruebas.

\begin{tablasimple}{>{\hsize=0.599\hsize\linewidth=\hsize}Y >{\hsize=0.804\hsize\linewidth=\hsize}Y >{\hsize=1.597\hsize\linewidth=\hsize}Y}
  {\cab{Rol} & \cab{Persona} & \cab{Responsabilidad en el protocolo}}
Contraparte Técnica & Designada por Distribuidora Puelche (BA, Art. 70.2, p. 37) & Revisa, observa y firma el acta \\
Jefe de Proyecto & Alex Aravena & Entrega el expediente y firma el acta por LafroX \\
Líder de Calidad y Pruebas & Maximiliano Miño & Verifica el expediente antes de la entrega; registra observaciones y su causa \\
Líder Funcional & Patricio Henríquez & Coordina los casos de aceptación con las personas del CLIENTE \\
Dueño del entregable & Responsable del paquete en el T-14 & Subsana las observaciones \\
\end{tablasimple}
\fuentesd{Fuente: (Bases Administrativas, art. 70.2, p. 37); (Subdocumento 1, sección 1.5); (Formulario T-13, sección 11).}

\section*{Referencias}
\addcontentsline{toc}{section}{Referencias}

Las fuentes siguientes sustentan los criterios citados en este documento; sus referencias se leen junto con el artículo, requisito o sección indicado en el lugar de uso.

\begin{itemize}[itemsep=0pt,parsep=0pt,topsep=2pt,before=\raggedright]
  \item Distribuidora Puelche S.A. (2026a). \emph{Bases Administrativas, Licitación TFEP-01/2026}.
  \item Distribuidora Puelche S.A. (2026b). \emph{Bases Técnicas Transversales, Licitación TFEP-01/2026}.
  \item Distribuidora Puelche S.A. (2026c). \emph{Bases Técnicas del Caso Logística}.
  \item Ley N.º 2.977. (1915). Diario Oficial de la República de Chile. Biblioteca del Congreso Nacional de Chile. \url{https://www.bcn.cl/leychile/navegar?idNorma=23639}
  \item Ley N.º 19.668. (2000). Diario Oficial de la República de Chile. Biblioteca del Congreso Nacional de Chile. \url{https://www.bcn.cl/leychile/navegar?idNorma=160270}
  \item Ley N.º 20.148. (2007). Diario Oficial de la República de Chile. Biblioteca del Congreso Nacional de Chile. \url{https://www.bcn.cl/leychile/navegar?idNorma=257080}
  \item Ley N.º 21.357. (2021). Diario Oficial de la República de Chile. Biblioteca del Congreso Nacional de Chile. \url{https://www.bcn.cl/leychile/navegar?idNorma=1161743}
  \item National Weather Service. (s.f.). \emph{Equinoxes and solstices through 2030} [Datos del U.S. Naval Observatory]. \url{https://www.weather.gov/dvn/Climate_Astronomical_Seasons}
  \item International Organization for Standardization. (2022). \emph{ISO/IEC/IEEE 42010:2022: Software, systems and enterprise—Architecture description}. \url{https://www.iso.org/standard/74393.html}
  \item Project Management Institute. (2017). \emph{Guía de los fundamentos para la dirección de proyectos (Guía del PMBOK)} (6.ª ed.). Project Management Institute.
\end{itemize}
